\documentclass{article}
\title{Evidence-Grounded Literature Assistance for LaTeX Papers}
\author{Yanban Evaluation Team}
\keywords{scientific writing, literature review, grounded generation}
\begin{document}
\maketitle

\section{Abstract}
This paper studies an assistant that helps authors revise LaTeX manuscripts while keeping citations and mathematical expressions intact. The assistant combines section-level polishing with evidence-grounded literature suggestions.

\section{Introduction}
Recent large language models can improve scientific writing fluency, but unconstrained generation may invent citations or alter technical notation. Retrieval-augmented methods reduce hallucination by grounding generated text in external documents\cite{lewis2020rag}. However, paper revision also requires preserving LaTeX structures such as Equation~\ref{eq:score} and figure references.

\section{Method}
For each section, protected LaTeX spans are masked before generation. The quality score is defined as
\begin{equation}
S = \alpha G + (1-\alpha) P,
\label{eq:score}
\end{equation}
where $G$ denotes grounding and $P$ denotes placeholder preservation.

\section{Discussion}
The current prototype focuses on text and bibliographic evidence. It does not yet evaluate figures or tables with a multimodal model.

\section{Conclusion}
A grounded revision assistant should improve readability without compromising traceability. Future work will add compile-time validation and human-in-the-loop patch acceptance.

\bibliographystyle{plain}
\bibliography{refs}
\end{document}
